\documentclass[10pt,reqno]{amsart}
\usepackage[T1]{fontenc}
\usepackage[utf8]{inputenc}
\usepackage{lmodern,amsmath,amssymb,mathtools,microtype,xurl}
\usepackage[a4paper,margin=25mm]{geometry}
\usepackage[hidelinks]{hyperref}
\newtheorem{theorem}{Theorem}
\newtheorem{lemma}[theorem]{Lemma}
\newtheorem{corollary}[theorem]{Corollary}
\title[Complete-shell collision bounds]{An explicit obstruction to principal domination\\in the Erd\H{o}s--Straus pair sum}
\author{The Clankers}
\date{September 19, 2026}
\begin{document}
\begin{abstract}
For the complete first-half divisor-pair source we prove that the usual
principal-minus-total-nonprincipal-magnitude estimate sums to at most
$-p\log p/24$ at every prime $p\equiv1\pmod4$, $p\ge2^{20}$. The same conclusion
holds when at most $J$ even-character modes are retained per shell once
$p\ge2^{20}J^3$. Exact integer collision bounds and proper even-character observations supply
original witnesses without assuming coset saturation.
A cubic-quotient example attains the integer bound while all principal-only
shell bounds at that prime are nonpositive. No universal ES lower bound is claimed.
\end{abstract}
\maketitle

\section{The original pair source}
Let $p\equiv1\pmod4$ be prime,
\[
 I_p=\{a\in\mathbb Z:p/4<a<p/2\},\quad R=4a-p,\quad
 c_a(g)=\#\{d\mid pa:d\equiv g\pmod R\}.
\]
All residues lie in $G=(\mathbb Z/R\mathbb Z)^\times$. Write
\[
 m_a=2\tau_{\rm div}(a),\quad \mathcal E_a=\sum_gc_a(g)^2,\quad
 T_a=\sum_gc_a(g)c_a(-g),\quad P_a=\frac{m_a^2}{\varphi(R)}.
\]
The canonical E/M criteria are $u\mid a^2$ and respectively $R\mid4u+1$ or
$R\mid4u+p$. Their normalized coordinates are
\[
 (h,r,s)=(d_0^2/u,u/d_0,a/d_0),\quad d_0=\gcd(a,u).
\]
The ordered returns are $(a,hs\kappa,phr\kappa)$ with
$R\kappa=pr+s$, or $(a,phs\lambda,phr\lambda)$ with $R\lambda=r+s$.
Clearing denominators verifies $4/p$, and the inverses are
$u=a^2/(Ry-pa)$ or $u=pa^2/(Ry-pa)$. These are the classical Type I/II
coordinates \cite[\S2]{ET}.

The original pair source is $b,c\mid pa$, $R\mid b+c$. Dividing by
$\gcd(b,c)$ returns an ordered M pair when the two $p$-valuations agree;
otherwise remove the one remaining $p$ and retain its side to obtain E.
If $h=a/(rs)$, the complete fibres are
\[
 M:(p^\epsilon dr,p^\epsilon ds),\quad d\mid h,\ \epsilon=0,1;
\quad
 E:(pdr,ds)\text{ or }(ds,pdr),\quad d\mid h.
\]
Their sizes are $2\tau_{\rm div}(h)$, so $T_a>0$ is exactly local occupancy.
The full first-half atlas is complete, and ES at $p$ is equivalent to
$\sum_{a\in I_p}T_a>0$; the expanded proof and all fibre maps are supplied
with the accompanying workbench source \cite{Workbench}.

Orthogonality gives
\[
 T_a=\frac1{\varphi(R)}\sum_\chi\chi(-1)
       \left|\sum_{d\mid pa}\chi(d)\right|^2,
 \qquad
 \mathcal E_a=\frac1{\varphi(R)}\sum_\chi
       \left|\sum_{d\mid pa}\chi(d)\right|^2.
\]
Thus $T_a\ge L_a:=2P_a-\mathcal E_a$. The following theorem addresses that
specific estimate, not every possible treatment of the signed sum.

\begin{theorem}\label{thm:neg}
For every prime $p\equiv1\pmod4$, $p\ge2^{20}$,
\[
 \boxed{\sum_{a\in I_p}L_a\le-\frac1{24}p\log p.}
\]
\end{theorem}
\begin{proof}
First $\tau_{\rm div}(n)<4n^{1/3}$. For primes $q\ge11$, $e+1\le q^{e/3}$;
for $q=2,3,5,7$ the maxima of $(e+1)/q^{e/3}$ occur at $e=3,2,1,1$ and
are at most $2,3/2,6/5,11/10$. Successive ratios prove the maxima, and the
product of these bounds is $3.96<4$.

Also
\[
 \sum_{\substack{r\le x\\r\equiv3\ (4)}}\frac1{\varphi(r)}
 <\frac75(1+\tfrac14\log x),\qquad x\ge3.
\]
Indeed $r/\varphi(r)=\sum_{d\mid r}\mu(d)^2/\varphi(d)$. Write $r=dm$;
the harmonic sum over the one required odd class of $m\pmod4$ is at most
$1+\frac14\log(x/d)$ by integration after its first term. Its coefficient sum is
\[
 \prod_{q\text{ odd prime}}(1+1/[q(q-1)])
 \le\exp\!\left(\sum_{j\ge1}\frac1{2j(2j+1)}\right)
 =e/2<7/5.
\]
All series converge absolutely. Consequently, using $a<p/2$ and
$2^{-2/3}<2/3$,
\[
 \mathcal P:=\sum_aP_a<\frac{224}{15}p^{2/3}(\log p+4).
\]
The original diagonal gives $\mathcal E_a\ge2\tau_{\rm div}(a)$. Put
$Y=\lfloor\sqrt p/2\rfloor$. Every $d\le Y$ dividing $a\in I_p$ has a distinct
complement above $Y$, and there are more than $p/(4d)-1$ such multiples of $d$.
Thus
\[
 \mathcal D:=\sum_a\mathcal E_a
 >p\sum_{d\le Y}\frac1d-4Y
 >p(\tfrac12\log p-\log2)-2\sqrt p.
\]
At the stated bound, $p^{1/3}>100$, $\log p>40/3$,
$\log2<3/4$, and $2/\sqrt p\le1/512$. Hence
\[
 \frac{2\mathcal P}{p\log p}<\frac{728}{1875},\qquad
 \frac{\mathcal D}{p\log p}>\frac7{16}.
\]
Since $7/16-728/1875>1/24$, the result follows.
\end{proof}

For a set $\Sigma_a$ of at most $J$ even characters, let
\[
 L_{a,\Sigma_a}=\frac2{\varphi(R)}\sum_{\chi\in\Sigma_a}
 \left|\sum_{d\mid pa}\chi(d)\right|^2-\mathcal E_a.
\]
Orthogonality gives $T_a\ge L_{a,\Sigma_a}\le2JP_a-\mathcal E_a$.
Thus the same summed negative bound holds at $p\ge2^{20}J^3$.
There is an effective all-sublinear extension. For $m\ge3$, define
\[
 K_m=\prod_{\substack{q<2^m\\q\ {\rm prime}}}
       \max_{0\le e<2m}\frac{(e+1)^m}{q^e},\qquad
 C_m=\min\{C\in\mathbb Z_{>0}:C^m\ge K_m\}.
\]
The consecutive ratio is below one once $e+1\ge2m$, since
$(1+1/(2m))^m<2$. Larger primes obey $e+1\le2^e$.
Hence $\tau_{\rm div}(a)\le C_ma^{1/m}$.
If $p\ge2^{20}$ and $p^{1-2/m}\ge12JC_m^2$, then
\[
 \frac{2J\mathcal P}{p\log p}<\frac{91}{300},\qquad
 \boxed{\sum_aL_{a,\Sigma_a}\le-\frac18p\log p.}
\]
Indeed $\mathcal P\le(7/5)C_m^2 2^{-2/m}p^{2/m}(\log p+4)$ and
$7/16-91/300>1/8$. At $m=4$ the finite product is
$K_4=127401984/25025$, so $C_4=9$ and
$p\ge\max\{2^{20},944784J^2\}$ suffices.
No zero-truncated or integer-refined bound is ruled out by these statements.

\section{An integer bound without a saturation assumption}
Let $B\le G$ contain $-1$, put $|B|=2s$, and let
$m_C=\sum_{g\in C}c_a(g)$ for each coset. This is a chosen observation
subgroup, not an asserted stabilizer of the actual support. For $m=sq+r$,
$0\le r<s$, set $\mathfrak b_s(m)=sq^2+r(2q+1)$.

\begin{theorem}\label{thm:int}
\[
 T_a\ge\mathcal L_{a,B}:=\sum_C\mathfrak b_s(m_C)-\mathcal E_a
 \ge L_{a,B}:=\frac2{|B|}\sum_Cm_C^2-\mathcal E_a.
\]
The first equality holds precisely when the antipodal-pair totals in every
coset differ by at most one. Also
$T_a\ge4\max(0,\lceil\mathcal L_{a,B}/4\rceil)$.
\end{theorem}
\begin{proof}
Partition each coset into $s$ pairs $\{g,-g\}$ and write $z_j=c_a(g)+c_a(-g)$.
Then $T_a+\mathcal E_a=\sum_C\sum_jz_j^2$. The least integer sum of squares
with fixed total $m_C$ is $\mathfrak b_s(m_C)$: moving one unit across a gap
of at least two decreases it. Real Cauchy--Schwarz proves the second bound.
Finally the commuting involutions $(b,c)\mapsto(c,b)$ and
$(b,c)\mapsto(pa/b,pa/c)$ act freely on the original target pairs.
The first cannot fix a pair because $R\nmid2b$; the second would make $pa$
a square; their product would give $bc=pa$ and make $-1$ a square modulo
$R\equiv3\pmod4$. Therefore $T_a$ is a multiple of four.
\end{proof}

The real bound retains all characters trivial on $B$, necessarily even, and
subtracts the exact remaining collision energy. It does not infer target
membership from a coarse occupied coset. If $[G:B]\le J$, then
$L_{a,B}\le2JP_a-\mathcal E_a$. The proof of Theorem~\ref{thm:neg} therefore gives
\[
 \boxed{\sum_aL_{a,B_a}\le-p\log p/24
 \quad\text{if }p\ge2^{20}J^3}
\]
for arbitrary adaptive per-shell choices of such subgroups. More generally the
same proof retains any at most $J$ even-character modes: each squared character
sum is at most $m_a^2$, and $T_a$ is twice the total even energy minus
$\mathcal E_a$. This statement does not apply to zero-truncating negative
bounds or to the integer refinement.

At the actual prime $p=87481$, complete enumeration of all 21870 first-half
shells gives $L_a\le0$ everywhere, even though 34 shells are occupied. The
integer refinement with $B=G$ is likewise nonpositive throughout. Thus no
nonnegative shell weighting of those bounds certifies this prime.

At $p=944329$, the principal-only bounds again all fail, but choose
\[
 a=236098=2\cdot97\cdot1217,
 \quad R=63,\quad B=\langle-1,13\rangle,
\]
\[
 B=\{1,8,13,20,22,29,34,41,43,50,55,62\}.
\]
The original mass is 16, its collision energy is 16, and the index-three coset
counts are $(8,8,0)$. Hence
\[
 \mathcal L_{a,B}=12+12-16=8=T_a,
 \qquad L_a=-16/9,
 \qquad L_{a,B}=16/3.
\]
With $\omega^2+\omega+1=0$, the retained nonprincipal coefficients are
$8(1+\omega)$ and its conjugate. Each has squared absolute value 64.
Both occupied coarse counts become zero modulo two; they retain their original
nonempty fibres. An actual pair is $(1,p\cdot1217)$, returning
\[
 (h,r,s,\kappa)=(194,1217,1,18242038),
\]
\[
 \frac4{944329}=\frac1{236098}+\frac1{3538955372}
                       +\frac1{4067138774169717196}.
\]
A separate occupied shell at this same prime has $R=47$, $T_a=60$ and negative
principal bound. Its group of order 46 has no proper subgroup containing $-1$
of order greater than two. The terminal order-two observation is the exact
target count, not a reduced universal certificate. A proper \emph{spectral}
observation does work there. At the original $a=236094$, take
$\chi(5)=\zeta_{23}$ modulo 47, so $\chi(-1)=1$. Form $C(X)$ by counting original
divisors of $pa$ at their discrete logarithms modulo 23. The Laurent norm
$Q(X)=C(X)C(X^{-1})\bmod(X^{23}-1)$ has
\[
 Q_0=192,\quad(Q_1,\ldots,Q_{11})=(187,176,159,138,116,92,70,50,33,21,14),
 \quad Q_{23-j}=Q_j.
\]
Exact rational Taylor evaluation gives
$1009<Q(\zeta_{23})<1010$ and $24<Q(\zeta_{23}^2)<25$.
The certificate uses $\pi=16\arctan(1/5)-4\arctan(1/239)$,
18 alternating terms for each arctangent and 18 cosine terms with their
next-term bounds; argument uncertainty is bounded using the unit Lipschitz
constant of cosine. The expanded proof gives the branch check and a second
independent rational enclosure. With collision energy 132 and principal square 2304,
retaining only $1,\chi^{\pm1},\chi^{\pm2}$ gives
\[
 T_a\ge\frac{2304+2Q(\zeta_{23})+2Q(\zeta_{23}^2)}{23}-132
 >58,\qquad T_a\ge60.
\]
Direct enumeration gives 60. These are five of 23 even modes, not the exact
terminal projection. The pair $(2,327)$ returns
$(236094,780326438241,4772638766)$.
The retained values and collision energy prove the bound; the omitted modes
are the projector's explicit kernel coordinates. Their favourable signs have
not been assumed.

\section{Reproduction and scope}
The standard-library package supplies complete original divisor/pair lists,
source factorizations, subgroup tables, inverse marks, rational constants, and
two code implementations. The small complete replay covers all 73798 shells
at $p\equiv1\pmod4$ primes through 3000. In the 4519 hard primes through two
million, the principal-only failures are exactly
$87481,196561,944329,1915201$; integer/coset certificates return witnesses at all
four. The separate checker exhausts every first-half shell of each failure.
No finite census proves a universal inequality. The analytic bound above is
proved separately, not inferred from those computations.

The prescribed-prime positivity of $\sum_aT_a$ remains unproved here. The
result rules out a specified absolute-value estimate and gives an exact
integer replacement retaining useful nonprincipal information. It does not
rule out signed arithmetic cancellation, and it makes no historical-priority,
Lean-build, or independent-review claim.

\begin{thebibliography}{2}
\bibitem{ET} Elsholtz, C., \& Tao, T. (2013).
Counting the number of solutions to the Erd\H{o}s--Straus equation on unit fractions.
\emph{Journal of the Australian Mathematical Society, 94}(1), 50--105.
\url{https://arxiv.org/abs/1107.1010v6}. Coordinates: \S2, Propositions 2.2 and 2.6.
\bibitem{Workbench} The Clankers. (2026, September 19).
\emph{Complete-shell collision bounds for the Erd\H{o}s--Straus equation}
[Turn 6 workbench and exact certificates]. Accompanying \texttt{core.tex},
\texttt{verify.py}, and \texttt{check\_independent.py}; exact antecedent pair
fibres retained from Turn 5, \S\S3--4.
\end{thebibliography}
\end{document}
